% $Header: /cvsroot/latex-beamer/latex-beamer/solutions/generic-talks/generic-ornate-15min-45min.en.tex,v 1.4 2004/10/07 20:53:08 tantau Exp $

%\documentclass[trans]{beamer}
%\documentclass[draft]{beamer}
\documentclass{beamer}

\usepackage{enumerate}
\usepackage{amscd,amsmath,xspace,amssymb,amsfonts,mathrsfs,latexsym}
\usepackage{graphicx}

\usepackage{array}
%\usepackage{mbboard}
%\usepackage{multicol}
%\usepackage[pdftex]{color}

%\includeonlyframes{current}

\usepackage[all,poly,knot,tips]{xy}

\newcommand{\lra}{\longrightarrow}
\newcommand{\lla}{\longleftarrow}
\newcommand{\Lra}{\Longrightarrow}
\newcommand{\Lla}{\Longleftarrow}

%% Script Shortcuts 
\def\CA{{\mathscr A}}
\def\CB{{\mathscr B}}
\def\CC{{\mathscr C}}
\def\CD{{\mathscr D}}
\def\CE{{\mathscr E}}
\def\CF{{\mathscr F}}
\def\CG{{\mathscr G}}
\def\CH{{\mathscr H}}
\def\CI{{\mathscr I}}
\def\CJ{{\mathscr J}}
\def\CK{{\mathscr K}}
\def\CL{{\mathscr L}}
\def\CM{{\mathscr M}}
\def\CN{{\mathscr N}}
\def\CO{{\mathscr O}}
\def\CP{{\mathscr P}}
\def\CQ{{\mathscr Q}}
\def\CR{{\mathscr R}}
\def\CS{{\mathscr S}}
\def\CT{{\mathscr T}}
\def\CU{{\mathscr U}}
\def\CV{{\mathscr V}}
\def\CW{{\mathscr W}}
\def\CX{{\mathscr X}}
\def\CY{{\mathscr Y}}
\def\CZ{{\mathscr Z}}

\newcommand{\A}{\ensuremath{\mathcal{A}}\xspace}
\newcommand{\B}{\ensuremath{\mathcal{B}}\xspace}
\newcommand{\Af}{\ensuremath{\mathcal{A}_f}\xspace}
\newcommand{\C}{\ensuremath{\mathcal{C}}\xspace}
\newcommand{\D}{\ensuremath{\mathcal{D}}\xspace}
\newcommand{\F}{\ensuremath{\mathcal{F}}\xspace}
\newcommand{\I}{\ensuremath{\mathcal{I}}\xspace}
\newcommand{\K}{\ensuremath{\mathcal{K}}\xspace}
\newcommand{\Kt}{\ensuremath{\mathcal{K}_t}\xspace}
\newcommand{\Kl}{\ensuremath{\mathcal{K}_{\ell}}\xspace}
\renewcommand{\L}{\ensuremath{\mathcal{L}}\xspace}
\renewcommand{\S}{\ensuremath{\mathcal{S}}\xspace}
\newcommand{\T}{\ensuremath{\mathcal{T}}\xspace}
\newcommand{\V}{\ensuremath{\mathcal{V}}\xspace}
\newcommand{\Vf}{\ensuremath{\mathcal{V}_f}\xspace}
\newcommand{\W}{\ensuremath{\mathcal{W}}\xspace}
\newcommand{\X}{\ensuremath{\mathcal{X}}\xspace}

\newcommand{\op}{\ensuremath{^{\textrm{op}}}\xspace}
\newcommand{\ps}{\ensuremath{^{\textrm{ps}}}\xspace}
\newcommand{\lax}{\ensuremath{^{\textrm{lax}}}\xspace}
\newcommand{\oplax}{\ensuremath{^{\textrm{oplax}}}\xspace}


\newcommand{\EMK}{\textnormal{\sf EM(\K)}\xspace}
\newcommand{\EM}{\textnormal{\sf EM}\xspace}
\newcommand{\EMwK}{\ensuremath{\EM^w(\K)}\xspace}

%\newcommand{\Cat}{\textnormal{\bf Cat}\xspace}
\newcommand{\Lex}{\textnormal{\sf Lex}\xspace}
\newcommand{\Lan}{\textnormal{\sf Lan}\xspace}
\newcommand{\Hom}{\textnormal{\sf Hom}\xspace}
\renewcommand{\t}{\times}
\newcommand{\ot}{\otimes}
\newcommand{\two}{\mathbf{2}\xspace}


\newtheorem{remark}[theorem]{Remark}
\newtheorem{proposition}[theorem]{Proposition}

\beamerdefaultoverlayspecification{<+->}

%\newenvironment{\next}{\uncover<+->{}{}}

\mode<presentation>
{
%  \usetheme[compress,subsection=false]{Szeged}
  % or ...
%  \useoutertheme[subsection=false]{miniframes}
%\useoutertheme[subsection=false, footline=authortitle, frame number]{miniframes}
\useoutertheme{infolines}
\setbeamercolor{separation line}{use=structure,bg=structure.fg!50!bg}
\usecolortheme{whale}
\usecolortheme{orchid}
\usecolortheme{rose}
\setbeamercolor{titlelike}{parent=structure}


%  \setbeamercovered{transparent}
  % or whatever (possibly just delete it)
}


\usepackage[english]{babel}
% or whatever

\usepackage[latin1]{inputenc}
% or whatever

%\usepackage{times}
%\usepackage[T1]{fontenc}
% Or whatever. Note that the encoding and the font should match. If T1
% does not look nice, try deleting the line with the fontenc.


\title[From lin alg to knots via categories] % (optional, use only with long paper titles)
{\textbf{From linear algebra to knot theory via categories}}

%{Presentation Subtitle} % (optional)

\author{Ross Street}
\institute{Macquarie University}
%[Universities of Somewhere and Elsewhere] % (optional, but mostly needed)
%{
%  \inst{1}%
%  Department of Computer Science\\
%  University of Somewhere
%  \and
%  \inst{2}%
%  Department of Theoretical Philosophy\\
%  University of Elsewhere}
% - Use the \inst command only if there are several affiliations.
% - Keep it simple, no one is interested in your street address.

\date[12 Aug 2013]{\small Mathematics Department Colloquium, E7B100} 

\subject{Colloquium}
% This is only inserted into the PDF information catalog. Can be left
% out. 

% Delete this, if you do not want the table of contents to pop up at
% the beginning of each subsection:
\AtBeginSubsection[]
{
  \begin{frame}<beamer>
    \frametitle{Outline}
    \tableofcontents[currentsection]
  \end{frame}
}
% If you wish to uncover everything in a step-wise fashion, uncomment
% the following command: 

\beamerdefaultoverlayspecification{<.->}

\begin{document}

%  \begin{frame}
%    \titlepage
%  \end{frame}


\begin{frame}%[plain]
%  \frametitle{titlepage}

  \begin{centering}
    {\usebeamerfont{title}\usebeamercolor[fg]{title}\inserttitle

\medskip 

{\small Ross Street}

{\small Macquarie University} 

\medskip
\uncover<1->{{}}
}

\vfill

\centerline{\usebeamerfont{title}\usebeamercolor[fg]{title}\insertdate}
% CT09, Cape Town}

  \end{centering}


\end{frame}


% \begin{frame}
%   \frametitle{Outline}
%   \tableofcontents
%   % You might wish to add the option [pausesections]
% \end{frame}

%\section[Stricts and weaks]{The strict/weak dichotomy}

%\subsection{Classical universal algebra}

\begin{frame}
  \frametitle{Purpose of the talk}

Monoidal categories are finding lots of uses in mathematics, physics and computer science. \\
Here I will be concerned with mathematical applications.
\pause

My goal is:
 \begin{itemize}
\item<1-> to recall some facts about vector spaces and linear functions \\ (MATH 136 and MATH 235);
\pause
\item<2-> to introduce categories and monoidal categories through the vector space example;
\pause
\item<3-> to define some concepts of monoidal category theory and identify them for vector spaces;
\pause
\item<4-> to explain string notation as a method for calculating;
\pause
\item<5-> to explain the two-way interaction between knot theory and categories.

\end{itemize}

\end{frame}

\begin{frame}
  \frametitle{Role of monoidal categories in mathematics}
  
  Simplistically:
  \bigskip
  
\[
\frac{\mathrm{Point \ \ in \ \ Euclidean \ \ Space}}{\mathrm{Vector \ \ Space}}=\frac{\mathrm{Vector \ \ Space}}{\mathrm{Monoidal \ \ Category}}\,.
\]
\end{frame}

\begin{frame}
  \frametitle{Left-hand side}

\noindent Formalizing properties of the Euclidean plane into the structure of vector space, 
we are able to transfer low-dimensional thinking to obtain precise results in higher dimensions.

\[
\xygraph{ !{\UseTips}
 !{/r6pc/: /u4pc/::}
 []*++[F-]+!C\txt{Vector spaces} !{+D*+!C\txt{ \ }} 
 ( [dl] *++[F-]+\txt{Plane\\geometry} : ?
 , ? : [dr] *++[F-]+\txt{Functional\\analysis}
 )
}\]
\end{frame}

\begin{frame}
  \frametitle{Right-hand side}
\noindent Formalizing properties of the totality of vector spaces into the structure of monoidal category, we are able to transfer ideas from linear algebra into some perhaps surprising areas.

\[
\xygraph{  !{\UseTips}!{/u4pc/::}
 !{/r6pc/: /u4pc/::}
 []*++[F-]+!C\txt{Monoidal categories} !{+D*+!C\txt{ \ }} 
 ( [dl] *++[F-]+\txt{Linear\\algebra} : ?
 , ? : [dr] *++[F-]+\txt{Knot\\theory}
 )
}\]
\end{frame}

\begin{frame}
 \frametitle{Origin of categories}
\begin{itemize}
\item Felix Klein (1872) recognized the importance of structure-respecting transformations in geometry.
They are what define the specific geometry in question.
\pause
\item These transformations were taken to be invertible, forming a groupoid.
\pause
\item General structure-respecting functions are also crucial.
\pause
\item The notion of ``category'' is a mathematical structure which captures these ideas.
\pause
\item The original defining paper of Samuel Eilenberg and Saunders Mac Lane (1945)
wanted to formalize what was meant by ``naturally'' or ``canonically'' isomorphic constructions in mathematics,
particularly in group theory.
\end{itemize}
\end{frame}

\begin{frame}
 \frametitle{The category of vector spaces}
\begin{itemize}
\item The \emph{category} $\mathrm{Vect}$ of (say, complex) vector spaces consists not only of the vector spaces $V$ (\emph{objects}) but also of the linear functions $f:V\lra W$ between them (\emph{morphisms}).
\pause
\item The basic operation is composition
 $$
 \xymatrix{ & W \ar[dr]^g &
                \\ V \ar[ur]^f  \ar[rr]_{g\circ f} & & X
                }
 $$ 
of the morphisms.
\pause
\item Each object $V$ has an identity morphism $1_V : V \lra V$.\end{itemize}
\end{frame}

\begin{frame}
\frametitle{Operations on vector spaces}

\begin{itemize}
\item A first course on linear algebra introduces the operation of \emph{direct sum} $V \oplus W$ of vector spaces and the quantity called \emph{dimension}. 
\pause
\item If $V$ and $W$ are finite dimensional,
\[
\dim (V \oplus W) = \dim V + \dim W.
\]
\pause
\item
\emph{Bilinear functions} $h : V \times W \lra X$ would also be defined.
\pause
\item
Yet the operation of \emph{tensor product}   $V \otimes W$ might await a course on multilinear algebra or module theory.
\pause
\item
For $V$ and $W$ finite dimensional, 
\[
\dim (V \otimes W) = \dim V \times \dim W.
\]
\pause
\item
Bilinear functions $h : V \times W \lra X$ are in canonical bijection with linear functions $f : V \otimes W \lra X$.
\end{itemize}
\end{frame}

\begin{frame}
\frametitle{Vector spaces of linear functions}
\begin{itemize}
\item There is also the vector space $X^V$ of linear functions from $V$ to $X$.
\pause
\bigskip
\item
Linear functions $f : V \otimes W  \lra X$ are in canonical bijection with linear functions
\[
g : W  \lra X^V.
\]
\end{itemize}
\end{frame}
\begin{frame}
\frametitle{Definition of category}
\begin{itemize}
\item 
A \emph{category}  $\CC$ consists of a set of \emph{objects} and, 
\pause
\item for each pair of objects $A, B$, a set $\CC (A,B)$ of morphisms $f : A  \lra B$, 
\pause
\item together with an associative composition rule $\circ$ ,
\pause
\item with identities $1_A : A  \lra A$,
\pause
\item subject to associativity and identity axioms.
\[  \xymatrix@R+2mm @C+2mm{
    A \ar[r]^{h \circ (g\circ f)} \ar[d]_f \ar[rd]^{g \circ f} & D   \\
    B \ar[r]_g    & C \ar[u]_h
  } \qquad
  \xymatrix@R+2mm @C+2mm{
    A \ar[r]^{(h \circ g)\circ f} \ar[d]_f  & D   \\
    B \ar[r]_g  \ar[ru]_{h \circ g}  & C \ar[u]_h}
 \qquad
   \xymatrix@R+2mm @C+2mm{
    A \ar[r]^f\ar[d]_{f=1_B \circ f}   & B \ar[ld]_{1_B}  \ar[d]^{g=g\circ 1_B}\\
    B \ar[r]_g     & C }
\]
\end{itemize}
\end{frame}

\begin{frame}
\frametitle{Commutative diagrams}
In a category, we can speak of \emph{commutative diagrams}:
\begin{equation*}
\xymatrix{
& Y \ar[rd]^-{ g}  & \\
X \ar[ru]^-{f} \ar[d]_-{u} \ar[rr]_-{w} & & Z \ar[d]^-{h} \\
V \ar[rr]_-{v} & & W }
\end{equation*}
%
\begin{center}
$h \circ w = v \circ u$ \quad and \quad $g \circ f = w$, \quad
so $v \circ u = h \circ g \circ f$.
\end{center}
\end{frame}

\begin{frame}
  \frametitle{Functors}

\begin{itemize}
\item Categories are themselves mathematical structures: so we should look at morphisms between them.
\pause
\bigskip
\item A \emph{functor} $T : \CC  \lra \CH$ assigns an object $TA$ of $\CH$ to each object $A$ of $\CC$, 
a morphism $Tf : TA  \lra TB$ in $\CH$ to each $f : A  \lra B$ in $\CC$, such that
\[
T1_A = 1_{TA} \quad \mbox{and} \quad T(g \circ f) = Tg \circ Tf.
\]
\end{itemize}
\end{frame}

\begin{frame}
  \frametitle{Examples of functors}
\begin{itemize}
\item 
Important constructions made on mathematical structures are functorial. 
\pause
\item Each vector space $V$ determines a functor
\[
T = V\otimes - : \mathrm{Vect}  \lra \mathrm{Vect} 
\]
defined by
\begin{align*}
TA &= V \otimes A  \\
Tf &= 1_V \otimes f : V \otimes A \lra V \otimes B\,,   \qquad  v \otimes a \longmapsto v \otimes f(a)\,. 
\end{align*}
\pause
\item In an obvious sense of \emph{product category}, tensor product is a functor
\[
- \otimes - : \mathrm{Vect}  \times \mathrm{Vect}  \lra \mathrm{Vect}\,.
\]
\end{itemize}
\end{frame}

\begin{frame}
  \frametitle{Other examples of functors}
\begin{itemize}
\item 
There is a functor $T_n : \mathrm{Vect}  \lra \mathrm{Vect} $
taking $V$ to the $n$-fold tensor power 
$$T_nV =  \underbrace{V\otimes V \otimes \dots \otimes V}_n = V^{\otimes n}\ .$$
\pause
\item Taking the direct sum of all these over $n$ gives
a functor
\[
T  : \mathrm{Vect}  \lra \mathrm{Vect} 
\]
defined by the `geometric series'
$$TV = \mathbb{C}\oplus V\oplus V^{\otimes 2}\oplus V^{\otimes 3}\dots \ .$$
\pause
\item In fact, each $TV$ is an algebra: indeed, we have a functor 
\[
T  : \mathrm{Vect}  \lra \mathrm{Alg} \ .
\]
\end{itemize}
\end{frame}

\begin{frame}
  \frametitle{Natural transformations}
\begin{itemize}
\item Categories exhibit a higher-dimensional structure in that there are morphisms between functors.
\pause
\item Suppose $S$ and $T : \CC \lra \CH$ are functors. A \emph{natural transformation}
\[
\theta : S \Longrightarrow T
\]
is a family of morphisms
\[
\theta_A : SA \lra TA
\]
in $\CH$, indexed by the objects $A$ of $\CC$, such that the square
\[
\xymatrix{
   S A \ar[r]^{\theta_A} \ar[d]_{Sf}   & TA  \ar[d]^{Tf}\\
    SB \ar[r]_{\theta_B}    & TB}
\]
commutes for all $f : A \lra B$ in $\CC$.
\end{itemize}
\end{frame}
\begin{frame}
  \frametitle{Examples of natural transformations}
\begin{itemize}
\item Each linear function $t : V \lra W$ determines a natural transformation 
$t \otimes - : V \otimes - \Lra W \otimes -$ with components
\begin{align*}
  t \otimes 1_A : &V \otimes A \lra W \otimes A \\
 & v \otimes a \lra t(v) \otimes a\,. 
\end{align*}
\pause
\item For each $n$, there is a natural transformation $s: T_n \Longrightarrow T_n$ 
whose component $s_V : V^{\otimes n}\lra V^{\otimes n}$ is defined by
$$s_V(v_1\otimes v_2 \otimes \dots \otimes v_n) = \frac{1}{n!} \sum_{\sigma} v_{\sigma(1)}\otimes v_{\sigma(2)} \otimes \dots \otimes v_{\sigma(n)}$$
where $\sigma$ runs over all the permutations of $n$. 
\pause
\item The componentwise image of $s$ gives the symmetric $n$-th-power functor $S_n : \mathrm{Vect}  \lra \mathrm{Vect}$.   
\end{itemize}
\end{frame}
\begin{frame}
  \frametitle{Monoidal category data}
A \emph{monoidal category} consists of a category $\CC$ equipped with a functor
\[
- \otimes - : \CC \times \CC \lra \CC,
\]
an object $I$, and invertible natural families of morphisms
\begin{align*}
 & \alpha^{}_{A,B,C} : (A \otimes B) \otimes C \lra A \otimes (B \otimes C)\\
 & \lambda_A : I \otimes A \lra A\,, \qquad \rho_A : A  \lra A\otimes I
\end{align*}
such that the following diagrams commute.
\end{frame}
\begin{frame}
  \frametitle{Monoidal category axioms}
\[
\xymatrix{
& (A \otimes (B \otimes C)) \otimes D \ar[rd]^-{\phantom{A}\alpha^{}_{A,B\otimes C,D}}  & \\
((A \otimes B )\otimes C) \otimes D \ar[ru]^-{\alpha^{}_{A,B,C}\otimes 1_{D}\phantom{A}} \ar[d]_-{\alpha^{}_{A\otimes B,C,D}} & & A \otimes ((B \otimes C) \otimes D ) \ar[d]^-{1_{A}\otimes\alpha^{}_{B,C,D}} \\
(A \otimes B) \otimes (C \otimes D) \ar[rr]_-{\alpha^{}_{A,B,C\otimes D }} & & A \otimes (B \otimes (C \otimes D)) 
}\] 
\[
\xymatrix{(A \otimes I)\otimes B \ar[rr]^{\alpha^{}_{A,I,B}}  
& & A \otimes (I \otimes B) \ar[d]^{1_A \otimes \lambda_B} \\
 A \otimes B \ar[rr]^{1_{A \otimes B}} \ar[u]^{\rho_A \otimes 1_B} & & A \otimes B}
\]

\begin{example} 
Take $\CC = \mathrm{Vect} $, usual tensor product of vector spaces as $\otimes$, and $I = \mathbb{C}$.
\end{example}
\end{frame}

\begin{frame}
  \frametitle{String diagrams}
\begin{itemize}
\item Penrose introduced \emph{string notation} to aid calculations in multilinear algebra during the 1950s (published 1971).
\pause
\item This was adapted to monoidal categories by Joyal-Street (published 1991).
\pause
\item In any category $\CC$, we can write $A \stackrel{f}{\lra} B$ or as the following. 

 \[ 
 \xygraph{{f} *\xycircle<8pt>{-}="m" "m"(-[u] [l(.3)d(.25)] {A},-[d] [l(.3)u(.25)] {B})}
\] 
\pause
\item Both notations are ``1-dimensional''.
\end{itemize}
\end{frame}
\begin{frame}
  \frametitle{String diagrams (continued)}
  
\begin{itemize}
\item In a monoidal category $\CC$, we use two dimensions: composition is up-down while tensoring is left-right.
\pause
\item Instead of $A \otimes C \stackrel{f}{\lra} B \otimes C \otimes A$, we draw.   
\[ 
\xygraph{{f} *\xycircle<8pt>{-}="m" "m"(-[l(.5)u] [l(.1)d(.4)] {A},-[r(.5)u] [r(.1)d(.4)] {C},-[l(.75)d] [l(.1)u(.4)] {B},-[d] [l(.15)u(.3)] {C},-[r(.75)d] [r(.1)u(.4)] {A})} 
\]
\pause
\item A morphism $I \stackrel{f}{\lra}  A \otimes B$ is drawn.
\[ 
\xygraph{!{0;(.5,0):(0,2)::}{f}*\xycircle<8pt>{-}="m" "m"(-[dl] [l(.15)u(.5)] {A},-[dr] [r(.15)u(.5)] {B})} 
\]
\end{itemize}
\end{frame}

\begin{frame}
  \frametitle{Progressive plane string diagrams}
\begin{itemize}
\item Notice that the monoidal structural items $\otimes, I, \alpha, \rho, \lambda$ do not appear.
\item
\[ 
\xygraph{ !{/r2.5pc/:}
{a} *\xycircle<6pt>{-}="a" [u(2)r(1.5)] {c} *\xycircle<6pt>{-}="c" [d(1)r(1.5)] {b} *\xycircle<6pt>{-}="b" [r(1.5)u(1.2)] {d} *\xycircle<6pt>{-}="d" "a" (-[u(3.2)]^-*{B},-[d]_-*{A},-"c"_-*{B}(-[u(1.2)]^-*{C},-"b"^-*{C}(-[d(2)]_-*{B},-"d"_-*{D}(-[u]_-*{D},-[d(3.2)]^-*{C}))))} 
\]
A diagram $\Gamma$ as above, labelled in $\CC$, is called a \emph{progressive plane string diagram}. 
It means that we have morphisms
\[
B \otimes B \stackrel{a}{\lra} A\,, \quad
C \otimes D \stackrel{b}{\lra} B\,,\quad
C   \stackrel{c}{\lra} B \otimes C\,,\quad
D \stackrel{d}{\lra}  D \otimes C\,.  
\]
\end{itemize}
\end{frame}

\begin{frame}
  \frametitle{The value of a string diagram}
\begin{itemize}
\item 
The \emph{value} of the diagram is the composite
\begin{align*}
 \mathrm{v}(\Gamma) =& (B \otimes C \otimes D \stackrel{1_B\otimes c \otimes d}\lra
 B \otimes B \otimes C \otimes D \otimes C
\\
 & \stackrel{1\otimes 1\otimes b \otimes 1}\lra   
 B \otimes B \otimes B \otimes C  \stackrel{a \otimes 1\otimes  1}\lra  
 A \otimes B  \otimes C)\,.
\end{align*}
\pause
\item It is obtained by composing the more elementary values of horizontal layers. 
\pause
\item This value is independent of \emph{deformation} in a natural sense. 
\pause
\item For example, the following diagram is a deformation of $\Gamma$ and has the same value as can be shown by the axioms for a monoidal category.   
\end{itemize}
\end{frame}
\begin{frame}

\[ 
\xygraph{{a} *\xycircle<6pt>{-}="a" [r(1.5)u] {c} *\xycircle<6pt>{-}="c" [r(1.5)d(2.5)] {b} *\xycircle<6pt>{-}="b" [r(1.5)u] {d} *\xycircle<6pt>{-}="d" "a" (-[u(2)]^-*{B},-[d(2.5)]_-*{A},-"c"_-*{B}(-[u]^-*{C},-"b"^-*{C}(-[d]_-*{B},-"d"_-*{D}(-[u(2.5)]_-*{D},-[d(2)]^-*{C}))))} 
\]
Its value is 
\begin{align*}
\mathrm{v}(\Gamma) =& (B \otimes C \otimes D \stackrel{1_B\otimes c \otimes 1}\lra 
 B \otimes B \otimes C \otimes D \stackrel{a \otimes 1 \otimes 1}\lra 
 A \otimes C \otimes D
\\
 & \stackrel{1\otimes 1\otimes d} \lra  
  A \otimes C \otimes D \otimes C   \stackrel{1 \otimes b \otimes  1} \lra  
  A \otimes B  \otimes C)\,.
\end{align*}
\end{frame}
\begin{frame}
The notation handles units well:
if $I   \stackrel{a}{\lra}   A \otimes B$ and $C \stackrel{b}{\lra}  I$, 
then the following three string diagrams all have the same value.
\[ 
\xygraph{!{0;(1.25,0):(0,.65)::} {\xybox{\xygraph{{\xybox{\xygraph{!{0;(.5,0):(0,2)::} {a} *\xycircle<6pt>{-}="m" "m"(-[dl] [l(.05)u(.5)] {A},-[dr] [r(.05)u(.5)] {B})}}} [u(.5)r(1.5)] {\xybox{\xygraph{{b} *\xycircle<6pt>{-}="m" "m" -[u] [r(.15)d(.25)] {C}}}}}}}
[dr] {,} [ur]
{\xybox{\xygraph{{\xybox{\xygraph{!{0;(.5,0):(0,1.24)::} {a} *\xycircle<6pt>{-}="m" "m"(-[dl] [l(.05)u(.5)] {A},-[dr] [r(.05)u(.5)] {B})}}} [u(.61)r(.41)] {\xybox{\xygraph{!{0;(1,0):(0,.62)::} {b} *\xycircle<6pt>{-}="m" "m" -[u] [r(.15)d(.25)] {C}}}}}}}
[dr] {,} [ur]
{\xybox{\xygraph{{\xybox{\xygraph{!{0;(.5,0):(0,2)::} {a} *\xycircle<6pt>{-}="m" "m"(-[dl] [l(.05)u(.5)] {A},-[dr] [r(.05)u(.5)] {B})}}} [u(.5)l(.5)] {\xybox{\xygraph{{b} *\xycircle<6pt>{-}="m" "m" -[u] [r(.15)d(.25)] {C}}}}}}}} 
\]
\end{frame}
\begin{frame}
\frametitle{Dual objects}
We can look at some concepts which make sense in any monoidal category.

A \emph{duality} between objects $A$ and $B$ is a pair of morphisms
\[
e : A \otimes B \lra I  \quad \mbox{and} \quad d : I \lra B \otimes A
\]
called the \textit{counit} and the \textit{unit}, respectively, satisfying

\[
 \xymatrix @R+3mm {
 A \otimes I \ar[r]^-{1\otimes d} \ar[rd]_{\rho_A^{-1}}^\cong 
 & A \otimes (B\otimes A) {\alpha^{-1} \atop {\mbox{\large $\cong$}}}  (A\otimes B)\otimes A \ar[r]^-{e\otimes 1}
 & I \otimes A \ar[ld]^{\lambda_B}_\cong \\
 & B }
\]
\[
 \xymatrix@R+3mm {
 I \otimes B \ar[r]^-{d\otimes 1} \ar[rd]_{\lambda_B}^\cong 
 & (B\otimes A) \otimes B {\alpha \atop {\mbox{\large $\cong$}}}  B\otimes (A\otimes B)\ar[r]^-{1\otimes e}
 & B \otimes I \ar[ld]^{\rho_B^{-1}}_\cong \\
 & B }
\]
\end{frame}
\begin{frame}
\frametitle{Duality axioms in string notation}
\[ 
\xygraph{!{0;(2.6,0):(0,.5)::} {\xybox{\xygraph{{e} *\xycircle{-}="e" [r(1.5)u] {d} *\xycircle{-}="a" "e"(-[u(2)l(.5)]^-*{A},-"a"_-*{B}-[d(2)]^(.75)*{A} [r(1.5)] -[u(3)]_(.16)*{A}) "a" [r(.75)d(.5)] {=}}}}  
 [r(1.0)] *\txt{,} [r(0.8)]
{\xybox{\xygraph{{d} *\xycircle{-}="a" [r(1.5)d] {e} *\xycircle{-}="e" "a"(-[d(2)l(.25)]^(.75)*{B},-"e"_-*{A}-[u(2)r(.5)]^-*{B} [r(1.5)] -[d(3)]^(.85)*{B}) "e" [r(1.125)u(.5)] {=}}}}} 
\]
Write $A \dashv B$  when such $e$ and $d$ exist: 

$A$ is a \emph{left dual} for $B$ while $B$ is a \emph{right dual} for $A$. 

\bigskip
Any two right duals for $A$ are isomorphic (if they exist). 

\bigskip
Write $A^\ast$ for a chosen right dual for~$A$. 
\end{frame}
\begin{frame}
Call $\CC$ right \emph{autonomous} (``compact'', ``rigid'') when each object has a right dual. 
Choice gives a functor
\[
( \ )^\ast : \CC^{\mathrm{op}}   \lra \CC
\]
where ``$\mathrm{op}$'' means $( \ )^\ast$ is ``contravariant'': 
$A \stackrel{f}{\lra} C$ goes to  $A^\ast \stackrel{\ f^\ast}{\longleftarrow} C^\ast $.

\[
\xymatrix{  & C\otimes C^\ast \ar[rd]^{e} \\
*+!L(.5){A \otimes C^\ast}\ar[ru]^{f\otimes 1} \ar[rd]_{1\otimes f^\ast} & & I \\
 & A\otimes A^\ast   \ar[ru]_{e} }  
\qquad\qquad
\xymatrix{  & A^\ast \otimes A \ar[rd]^-{1\otimes f} \\
I \ar[ru]^-{d} \ar[rd]_-{d} & & *+!R(.5){A^\ast \otimes C} \\
 & C^\ast \otimes C   \ar[ru]_-{f^\ast \otimes 1} }
 \]
\end{frame}
\begin{frame}
We suppress $e$ and $d$ from notation by writing them as
%
\[ \kern-1cm
\xygraph{!{0;(-2.5,0):(0,1)::} [dd]*{}="l" [r] *{}="r" 
 "l"-@`{"l"+(.2,1),"r"+(-.2,1)}"r"|*++!U{e} 
"l" [l(.125)u(.35)] {A^*} "r" [r(.125)u(.35)] {A} 
}
\kern-3cm
\xygraph{!{0;(2.5,0):(0,1)::} [d] *{}="l" [r] *{}="r"
 "l"-@`{"l"+(.2,1),"r"+(-.2,1)}"r" |*++!D{d}
"l" [l(.125)u(.35)] {A^*} "r" [r(.125)u(.35)] {A}  
} 
\]
\end{frame}
\begin{frame}
The duality conditions become string straightening rules.
\[
\xygraph{ !{/r1.4pc/: /u2.75pc/::}
 []*+!D{A}
!{\xcapv-@(0)}
!{\vcap[-]}
[r] (?*!DR{A^*} , !{\vcap} )
[r] !{\xcapv-@(0)} *+!U{A}
} \quad = 
\xygraph{
!{/r1.3pc/: /u2pc/::}
!{\xcapv-@(0)}
!{\xcapv-@(0)} *+!U{A}
} \qquad,\qquad
\xygraph{ !{/r1.4pc/: /u2.75pc/::}
 [rr]*+!D{A^*}
  !{\xcapv-@(0)}
 [l] !{\vcap[-]}
 (?*!DL{A} ,  [l] !{\vcap} )
 [l] !{\xcapv-@(0)} *+!U{A^*}
} \quad = 
\xygraph{
!{/r1.3pc/: /u2pc/::}
!{\xcapv-@(0)}
!{\xcapv-@(0)} *+!U{A^*}
} \]

\end{frame}
\begin{frame}
This allows anticlockwise looping by taking $( \ )^\ast$.

$f^\ast$ is depicted as follows. 
\[ 
f^* \;\; : 
\xygraph{ !{/r2pc/: /u2pc/::}
 [rr]*+!D{B^*}
  !{\xcapv-@(0)} !{\xcapv-@(0)}
 [l] !{\vcap[-]} - [u]*+<8pt>[Fo]{f} _(.3)*{B}
 ( ?-[u]_(.8)*{A} [l] !{\vcap} )
 [lu] !{\xcapv-@(0)\xcapv-@(0)} *+!U{A^*}
} \]
%

Diagrams are not necessarily ``progressive''  any more.   

They can backtrack.
\end{frame}
\begin{frame}
\begin{example}
An object of $\mathrm{Vect} $ has a right (left) dual if and only if it is finite dimensional (exercise!). 
\bigskip

If $V$ has basis $v_1,\ldots,v_n$ and $v^\ast_1,\ldots,v^\ast_n$ is the ``dual basis'' of $V^\ast = \mathbb{C}^V$ then we have

\begin{align*}
 V \otimes V^\ast &\stackrel{e}{\lra} \mathbb{C}    
   & \mathbb{C}
     &\stackrel{d}{\lra}   V^\ast \otimes V  \\
  v \otimes \varphi &\longmapsto \varphi (v)   
   & 1 &\longmapsto \sum^n_{i=1} v^\ast_i \otimes v_i\,.
\end{align*}
\end{example}
\end{frame}
\begin{frame}
\frametitle{Braidings}
A \textit{braiding} on a monoidal category $\CC$ is a natural family of invertible morphisms 
\[
\gamma^{}_{A,B} : A \otimes B  \stackrel{\mbox{{\small  $\cong$}}}{\lra} 
 B \otimes A
\]
satisfying (ignoring $\alpha$)
\[
 \xymatrix @C-10mm{
A \otimes B \otimes C \ar[rr]^{\gamma_{A\otimes B,C}^{}} \ar[rd]_-{1\otimes \gamma_{B,C}^{}} 
& & C \otimes A \otimes B  \\
& A \otimes C \otimes B\ar[ru]_-{\gamma_{A,C}^{} \otimes 1}}  
 \qquad 
\xymatrix @C-10mm{
A \otimes B \otimes C \ar[rr]^{\gamma_{A, B\otimes C}^{}} \ar[rd]_-{\gamma_{A,B}^{} \otimes 1} 
& & B \otimes C \otimes A  \\
& B \otimes A \otimes C \ar[ru]_-{1\otimes \gamma_{A,C}^{}}}   
\]

\begin{example}
In $\mathrm{Vect}$, $\gamma^{}_{A,B} (a\otimes b) = b \otimes a$.
\end{example}
\end{frame}
\begin{frame}

\frametitle{Enter three dimensions!}
By writing $\gamma$ as a \emph{crossing}
\[
\xygraph{
 []{\gamma^{}_{A,B} :} [r][u(0.5)] *+!DR{A}( !{\xunderv}, [r]*+!DL{B} , [dr]*+!UL{A} ,  [d]*+!UR{B} )}
\]
our string diagrams move into three-dimensional Euclidean space.   

The axioms say
\[
\xygraph{[u(.75)]*+{A} ( [r]*+{B} , [d(1.5)l(.5)]*+{B'} , [dr(1.5)]*+{A'},
  [d]*++=[o][Fo]{b}="b"-"B'" , [dr]*++=[o][Fo]{a}="a"-"A'" ,
!{\xunderv ~{"A"}{"B"}{"b"}{"a"}} 
  )}  = 
\xygraph{[u(.75)]*+{A} ( [rr]*+{B} , [d(1.5)r(.5)]*+{B'} , [dr(1.5)]*+{A'},
  [dr(.5)r]*++=[o][Fo]{b}="b"-"B" , [dr(.5)]*++=[o][Fo]{a}="a"-"A" ,
!{\xunderv ~{"a"}{"b"}{"B'"}{"A'"}} 
  )}  
\]
\end{frame}
\begin{frame}
\[
\xygraph{ !{\knotholesize{15pt}}
!{/r1.8pc/: /u4pc/::}
[d(.5)] ( ?*+!DR{A} , [r] *+!DL{C}
 , !{\xunderv @(.45)}
, [r(.3)] ( ?*+!D{B} , ?- [rd]
))
}
=
\xygraph{
!{/r1.3pc/: /u2pc/::}
%( ?*+!DR{A} , [r] *+!D{B} [r] *+!DL{C} ,
 !{\xcapv-@(0)}
!{\vtwistneg} ( []*+!U{C}  [r]*+!U{A}  [r]*+!U{B} ,
[uur] !{\vtwistneg}
[r]!{\xcapv@(0)}
)
}
\quad,\qquad
\xygraph{ %!{\knotholesize{20pt}}
!{/r1.8pc/: /u4pc/::}
[d(.5)] ( ?*+!DR{A} , [r(.7)] *+!D{B}  , [r] *+!DL{C} ,
!{\xunderv}
, [r(.7)] - [ld] |(.35)*{\khole}
}
= 
\xygraph{
!{/r1.3pc/: /u2pc/::}
[r] !{\vtwistneg}
!{\xcapv-@(0)}
[uurr]
!{\xcapv@(0)}
[l] !{\vtwistneg} 
[l]*+!U{B}  [r]*+!U{C}  [r]*+!U{A} 
}\,.
\]
\bigskip

A monoidal category with chosen braiding is called \emph{braided}.
\end{frame}
\begin{frame}
The braiding automatically satisfies what is called the \textit{Yang--Baxter equation} 
or the \textit{braid identity} or a \textit{Reidemeister move} in different contexts. 
Here is the proof:

\[
\xymatrix{
  & B\otimes A \otimes C \ar[r]^-{1\otimes \gamma} 
    & B\otimes C \otimes A \ar[rd]^-{\gamma \otimes 1} \\
 A\otimes B\otimes C 
   \ar[ru]^-{\gamma \otimes 1} \ar[rru]_-{\;\gamma^{}_{A,B\otimes C}} 
    \ar[rd]_-{1\otimes \gamma} \ar@{}[rrr]|{naturality} 
    & & & C\otimes B \otimes A \\
 & A\otimes C \otimes B \ar[rru]^-{\gamma^{}_{A,C\otimes B}} \ar[r]_-{\gamma \otimes 1} 
   & C\otimes A \otimes B \ar[ru]_-{1\otimes \gamma} }
\]
\kern-20pt
\[
\xygraph{
!{/r1.3pc/: /u2pc/::}
!{\vtwistneg}
!{\xcapv-@(0)}
!{\vtwistneg}
[uuurr]
!{\xcapv@(0)}
[l]!{\vtwistneg}
[r]!{\xcapv@(0)}
[r][u(1.5)]
{=}
[r][u(1.5)]
!{\xcapv-@(0)}
!{\vtwistneg}
!{\xcapv-@(0)}
[uuur]
!{\vtwistneg}
[r]!{\xcapv@(0)}
[l]!{\vtwistneg}
}
\]
\end{frame}
\begin{frame}
\frametitle{The braid category $\mathbb{B}$}
Here is an example of a braided monoidal category which is vastly different from $\mathrm{Vect}$.
Objects are natural numbers 0, 1, 2, \ldots
\[
\mathbb{B}(m,n) = \begin{cases}
\emptyset  &  \text{when } m \neq n  \\
\mbox{braid group} \ B_n \text{ on } n \text{ strings } & \text{when } m=n
\end{cases}
\]
The functor $\mathbb{B} \times \mathbb{B} \stackrel{\otimes}{\lra} \mathbb{B}$ is defined on objects by 
$(m,n) \longmapsto m+n$ and on morphisms as indicated by:
\[
\xy
0;/r.18pc/:
(0,0)*{\xy
(-15,20);(25,20) **\dir{-};
(-15,0);(25,0) **\dir{-};
(-10,20)*{\bullet}="T1"; 
(0,20)*{\bullet}="T2"; 
(10,20)*{\bullet}="T3"; 
(20,20)*{\bullet}="T4";
(-10,0)*{\bullet}="B1"; 
(0,0)*{\bullet}="B2"; 
(10,0)*{\bullet}="B3"; 
(20,0)*{\bullet}="B4";
(-5,15)*{\hole}="a";
(-5,15)*{}="a'";
(-5,8)*{\hole}="b";
(-5,8)*{}="b'";
"T1";"a'" **\crv{"T1"+(0,-3)&"a'"+(-2,2)};
"T2";"a" **\crv{"T2"+(0,-3)&"a"+(2,2)};
"a'";"b" **\crv{"a'"+(2,-2)&"b"+(2,2)};
"a";"b'" **\crv{"a"+(-2,-2)&"b'"+(-2,2)};
"b";"B1" **\crv{"b"+(-2,-2)&"B1"+(0,5)};
(15,12)*{\hole}="c";
(15,12)*{}="c'";
"T3";"c'" **\crv{"T3"+(0,-3)&"c'"+(-2,2)};
"T4";"c" **\crv{"T4"+(0,-3)&"c"+(2,2)};
"c";"B2" **\crv{"c"+(2,-2)&"B2"+(0,3)} \POS?(.5)*{\hole}="d";
"b'";"d" **\crv{"b'"+(2,-2)&"d"+(-2,1)};
"d";"B4" **\crv{"d"+(2,-1)&"B4"+(0,3)} \POS?(.5)*{\hole}="e";
"c'";"e" **\crv{"c'"+(2,-2)&"e"+(1,2)};
"e";"B3" **\crv{"e"+(-1,-2)&"B3"+(0,3)};
\endxy};
(25,0)*{\otimes};
(40,0)*{\xy
(-15,20);(5,20) **\dir{-};
(-15,0);(5,0) **\dir{-};
(-10,20)*{\bullet}="T1"; 
(0,20)*{\bullet}="T2"; 
(-10,0)*{\bullet}="B1"; 
(0,0)*{\bullet}="B2"; 
(-5,14)*{\hole}="a";
(-5,14)*{}="a'";
(-5,6)*{\hole}="b";
(-5,6)*{}="b'";
"T1";"a" **\crv{"T1"+(0,-3)&"a"+(-2,2)};
"T2";"a'" **\crv{"T2"+(0,-3)&"a'"+(2,2)};
"a";"b'" **\crv{"a"+(2,-2)&"b'"+(2,2)};
"a'";"b" **\crv{"a'"+(-2,-2)&"b"+(-2,2)};
"b'";"B1" **\crv{"b'"+(-2,-2)&"B1"+(0,5)};
"b";"B2" **\crv{"b"+(2,-2)&"B2"+(0,5)};
\endxy};
(55,0)*{=};
(90,0)*{\xy
(-15,20);(45,20) **\dir{-};
(-15,0);(45,0) **\dir{-};
(-10,20)*{\bullet}="T1"; 
(0,20)*{\bullet}="T2"; 
(10,20)*{\bullet}="T3"; 
(20,20)*{\bullet}="T4";
(30,20)*{\bullet}="T5";
(40,20)*{\bullet}="T6";
(-10,0)*{\bullet}="B1"; 
(0,0)*{\bullet}="B2"; 
(10,0)*{\bullet}="B3"; 
(20,0)*{\bullet}="B4";
(30,0)*{\bullet}="B5";
(40,0)*{\bullet}="B6";
(-5,15)*{\hole}="a";
(-5,15)*{}="a'";
(-5,8)*{\hole}="b";
(-5,8)*{}="b'";
"T1";"a'" **\crv{"T1"+(0,-3)&"a'"+(-2,2)};
"T2";"a" **\crv{"T2"+(0,-3)&"a"+(2,2)};
"a'";"b" **\crv{"a'"+(2,-2)&"b"+(2,2)};
"a";"b'" **\crv{"a"+(-2,-2)&"b'"+(-2,2)};
"b";"B1" **\crv{"b"+(-2,-2)&"B1"+(0,5)};
(15,12)*{\hole}="c";
(15,12)*{}="c'";
"T3";"c'" **\crv{"T3"+(0,-3)&"c'"+(-2,2)};
"T4";"c" **\crv{"T4"+(0,-3)&"c"+(2,2)};
"c";"B2" **\crv{"c"+(2,-2)&"B2"+(0,3)} \POS?(.5)*{\hole}="d";
"b'";"d" **\crv{"b'"+(2,-2)&"d"+(-2,1)};
"d";"B4" **\crv{"d"+(2,-1)&"B4"+(0,3)} \POS?(.5)*{\hole}="e";
"c'";"e" **\crv{"c'"+(2,-2)&"e"+(1,2)};
"e";"B3" **\crv{"e"+(-1,-2)&"B3"+(0,3)};
(35,14)*{\hole}="g";
(35,14)*{}="g'";
(35,6)*{\hole}="h";
(35,6)*{}="h'";
"T5";"g" **\crv{"T5"+(0,-3)&"g"+(-2,2)};
"T6";"g'" **\crv{"T6"+(0,-3)&"g'"+(2,2)};
"g";"h'" **\crv{"g"+(2,-2)&"h'"+(2,2)};
"g'";"h" **\crv{"g'"+(-2,-2)&"h"+(-2,2)};
"h'";"B5" **\crv{"h'"+(-2,-2)&"B5"+(0,5)};
"h";"B6" **\crv{"h"+(2,-2)&"B6"+(0,5)};
\endxy};
\endxy
\]
The isomorphisms $\alpha, \lambda, \rho$ are identities.

\end{frame}
\begin{frame}
The braiding $\gamma^{}_{m,n} $ is indicated by:
\[
\xy
(-25,18)*+{4+2};(-25,0)*+{2+4} **\dir{-} \POS?(.5)*+!R{\gamma_{4,2}} ?>*\dir{>};
(5,22)*+{m};(-10,22)*{} **\dir{-} ?>*\dir{>};
(5,22)*+{\hole};(20,22)*{} **\dir{-} ?>*\dir{>};
(35,22)*+{n};(30,22)*{} **\dir{-} ?>*\dir{>};
(35,22)*+{\hole};(40,22)*{} **\dir{-} ?>*\dir{>};
(-15,18);(45,18) **\dir{-};
(-15,0);(45,0) **\dir{-};
(-10,18)*{\bullet}="T1"; 
(0,18)*{\bullet}="T2"; 
(10,18)*{\bullet}="T3"; 
(20,18)*{\bullet}="T4";
(30,18)*{\bullet}="T5";
(40,18)*{\bullet}="T6";
(-10,0)*{\bullet}="B1"; 
(0,0)*{\bullet}="B2"; 
(10,0)*{\bullet}="B3"; 
(20,0)*{\bullet}="B4";
(30,0)*{\bullet}="B5";
(40,0)*{\bullet}="B6";
"T1";"B3" **\crv{} \POS?(.8333)*{\hole}="h" \POS?(.8)*{\hole}="d";
"T2";"B4" **\crv{} \POS?(.6667)*{\hole}="g" \POS?(.75)*{\hole}="c";
"T3";"B5" **\crv{} \POS?(.5)*{\hole}="f" \POS?(.6667)*{\hole}="b";
"T4";"B6" **\crv{} \POS?(.3333)*{\hole}="e" \POS?(.5)*{\hole}="a" ;
"T5";"a" **\crv{};
"a";"b" **\crv{};
"b";"c" **\crv{};
"c";"d" **\crv{};
"d";"B1" **\crv{};
"T6";"e" **\crv{};
"e";"f" **\crv{};
"f";"g" **\crv{};
"g";"h" **\crv{};
"h";"B2" **\crv{};
\endxy
\]

\begin{proposition} 
$\mathbb{B}$ is the free braided monoidal category generated by a single object.
\end{proposition}
\end{frame}
\begin{frame}
The single generating object is $1$ and tensoring the braiding $\gamma^{}_{1,1}$ on both sides with identity morphisms gives the generators $s_i$ for the braid group $B_n$ as indicated:
\[
\xy
(-30,18)*+{n};(-30,0)*+{n} **\dir{-} \POS?(.5)*+!R{s_i} ?>*\dir{>};
(-19,18);(-9,18) **\dir{-};
(-9,18);(-5,18) **\dir{.};
(-5,18);(35,18) **\dir{-};
(35,18);(39,18) **\dir{.};
(39,18);(49,18) **\dir{-};
(-19,0);(-9,0) **\dir{-};
(-9,0);(-5,0) **\dir{.};
(-5,0);(35,0) **\dir{-};
(35,0);(39,0) **\dir{.};
(39,0);(49,0) **\dir{-};
(-14,18)*{\bullet}="T1"+(0,4)*{1};
(0,18)*{\bullet}="T2"+(0,4)*{i-1};
(10,18)*{\bullet}="T3"+(0,4)*{i};
(20,18)*{\bullet}="T4"+(0,4)*{i+1};
(30,18)*{\bullet}="T5"+(0,4)*{i+2};
(44,18)*{\bullet}="T6"+(0,4)*{n};
(-14,0)*{\bullet}="B1"; 
(0,0)*{\bullet}="B2"; 
(10,0)*{\bullet}="B3"; 
(20,0)*{\bullet}="B4";
(30,0)*{\bullet}="B5";
(44,0)*{\bullet}="B6";
"T1";"B1" **\dir{-};
"T2";"B2" **\dir{-};
(15,9)*{\hole}="a";
(15,9)*{}="a'";
"T3";"a'" **\crv{"T3"+(0,-4)&"a'"+(-2,2)};
"T4";"a" **\crv{"T4"+(0,-4)&"a"+(2,2)};
"a";"B3" **\crv{"a"+(-2,-2)&"B3"+(0,4)};
"a'";"B4" **\crv{"a'"+(2,-2)&"B4"+(0,4)};
"T5";"B5" **\dir{-};
"T6";"B6" **\dir{-};
\endxy
\]
The Proposition means that there is an equivalence of categories
\[
\mathrm{BStMon}(\mathbb{B},\CC) \simeq \CC
\]
between the category of braiding-and-tensor-preserving functors $\mathbb{B} \to \CC$ and the category $\CC$ itself, for every braided monoidal category $\CC$.
\end{frame}

\begin{frame}
\frametitle{Yang-Baxter operators}
A \emph{Yang--Baxter operator} $y : A \otimes A \to A \otimes A$ on an object  $A$ of any monoidal category $\CC$ is an invertible morphism  satisfying commutativity of the following hexagon. There is an obvious category $\mathrm{YB}\CC$ of such pairs $(A,y)$.

\begin{equation*}
\xymatrix{& A^{\otimes 3} \ar[r]^{1\otimes y} & A^{\otimes 3} \ar[rd]^{y\otimes 1} \\
 A^{\otimes 3}\ar[ru]^{y\otimes 1}  \ar[rd]_{1\otimes y}  & & & A^{\otimes 3} \\
 & A^{\otimes 3}  \ar[r]_{y\otimes 1 \ } & A^{\otimes 3} \ar[ru]_{1\otimes y} }
 \end{equation*}
If $\CC$ is braided, each object $A$ has a Yang--Baxter operator on it, namely,
\[
y =  \gamma^{}_{A,A} : A \otimes A \lra A \otimes A\,.
\]
\end{frame}

\begin{frame}

\begin{proposition} 
$\mathbb{B}$  is the free monoidal category generated by an object bearing a Yang--Baxter operator.   That is, for every monoidal category $\CC$, 
\[
\mathrm{StMon}(\mathbb{B},\CC) \simeq \mathrm{YB}\CC\,.
\]
\end{proposition} 
\begin{proposition} 
Any right autonomous braided monoidal category is left autonomous.
\end{proposition} 

To see this we take $A^\ast \otimes A \stackrel{\gamma}{\lra}  A \otimes A^\ast   \stackrel{e}{\lra}  I$ 
and $I   \stackrel{d}{\lra}  A^\ast \otimes A \stackrel{\gamma^{-1}}{\lra}  A \otimes A^\ast$ as counit 
and unit for $A^\ast$ as left dual to $A$. 
\end{frame}

\begin{frame}
\frametitle{Trace of an endomorphism}
The \emph{trace} of an endomorphism
$ f : A    \lra   A $
in an autonomous braided monoidal category is defined by
\[
\mathrm{Tr} (f) = \Bigl(I \stackrel{d}{\lra} A^\ast \otimes A \stackrel{1 \otimes f}{\lra}   
  A^\ast \otimes A \stackrel{e'}{\lra}  I\Bigr)\,.
\]
\[
\xygraph{
!{/r1.3pc/: /u1.6pc/::}
!{\vcap} ( [r] - [dd] ^(.3)*{A}|(.8)*+[Fo]{f}
, ? - [dd] _(.3)*{A^*}
!{\vtwistneg}
%!{\xcapv-@(0)}
!{\vcap[-]}
)}
\]
\end{frame}
\begin{frame}
\begin{example} 
If $V \in \mathrm{Vect} $ is finite dimensional with basis $v^{}_1,\ldots,v^{}_n$ and
$f$ is defined by $f(v^{}_i) = \sum_j a^{}_{ij} v^{}_j$, then
\begin{align*}
\mathrm{Tr}(f) :  \mathbb{C}  &\lra \mathbb{C} \\
                  1 &\longmapsto \displaystyle \sum_i a^{}_{ii}\,.
\end{align*}
\end{example}

\end{frame}
\begin{frame}
\frametitle{The category $\mathbb{T}$ of tangles }
This is a monoidal category first defined by Yetter (1988). 

The objects are words $- + + - - +$ in symbols $+$ and $-$.    

Morphisms are tangles
\[
\xy
(0,20)*+!D{-}="t1";(10,20)*+!D{-}="t2";(20,20)*+!D{-}="t3";
(30,20)*+!D{+}="t4";(40,20)*+!D{+}="t5"; (50,20)*+!D{+}="t6";(60,20)*+!D{-}="t7";
(0,-10)*+!U{-}="b1";(10,-10)*+!U{+}="b2";(20,-10)*+!U{-}="b3";
"t3"+(3,-8)="a";
"t2"+(8,-15)="c";
"t3";"a"*{\hole} **\crv{"t3"+(0,-3)&"a"+(-1,2)} \POS?(.5)="b";
"t4";"b"*{\hole} **\crv{"t4"+(0,-1)&"b"+(1,1)} \POS?(.7)*\dir{>};
"t2";"c"*{\hole} **\crv{"t2"+(0,-2)&"c"+(-3,2)} \POS?(.8)="d" \POS?(.6)="e";
"b"*{\hole};"d"*{\hole} **\crv{"b"-(2,2)&"d"+(3,1)} \POS?(.5)="f";
"e"*{\hole};"f"*{\hole} **\crv{"f"};
"f"*{\hole};"a" **\crv{"f"&"a"+(-2,0)};
"e"+(-8,-2)="g";
"e"*{\hole};"g" **\crv{"e"+(1,0)&"g"+(2,1)} \POS?(.5)*\dir{<};
"t1";"g"*{\hole} **\crv{"t1"+(0,-3)&"g"+(0,3)} \POS?(.5)*\dir{<};
"d"+(-4,-4)="h";
"h"*{\hole};"b1" **\crv{"h"+(-4,0)&"b1"+(1,2)} \POS?(.5)="j" \POS?(.25)="k" \POS?(.65)*\dir{<};
"h"*{\hole};"c" **\crv{"h"+(2,0)&"c"+(-2,0)};
"d"*{\hole};"h" **\crv{"d"-(3,1)&"h"+(0,2)};
"g";"g"+(-5,-3)="i" **\crv{"g"+(-2,-1)&"i"+(2,2)};
"i";"i"*{\hole} **\crv{"i"-(8,8)&"i"+(-8,8)};
"i"*{\hole};"j"*{\hole} **\crv{"i"+(2,-2)&"j"+(-2,2)};
"g"*{\hole};"k"*{\hole} **\crv{"g"-(0,2)&"k"+(-3,1)};
"k"*{\hole};"b3" **\crv{"k"+(3,-1)&"b3"+(0,3)} \POS?(.3)="m" \POS?(.6)="l";
"j"*{\hole};"l"*{\hole} **\crv{"j"+(2,-2)&"l"+(-3,0)} \POS?(.4)="n";
"h";"m"*{\hole} **\crv{"h"-(0,2)&"m"+(2,2)};
"m"*{\hole};"n"*{\hole} **\crv{"m"-(2,2)&"n"+(1,2)};
"n"*{\hole};"b2" **\crv{"n"-(1,2)&"b2"+(0,2)};
"l"+(12,-1)="o";
"l"*{\hole};"o"*{\hole} **\crv{"l"+(2,0)&"o"+(-2,0)};
"c"*{\hole};"o" **\crv{"c"-(-3,2)&"o"+(-1,3)};
"o";"t6" **\crv{"o"+(4,-12)&"t6"-(0,5)} \POS?(.6)*\dir{<};
"a"+(6,-4)="p";
"a";"p"*{\hole} **\crv{"a"+(3,0)&"p"+(-2,2)};
"p"*{\hole};"o"*{\hole} **\crv{"p"+(2,-2)&"o"+(4,0)} \POS?(.5)="q";
"a"*{\hole};"q"*{\hole} **\crv{"a"+(1,-2)&"q"+(-4,0)} \POS?(.45)="r";
"q"*{\hole};"p" **\crv{"q"+(6,0)&"p"+(4,2)};
"p";"r"*{\hole} **\crv{"p"-(2,1)&"r"+(2,1)};
"r"*{\hole};"c" **\crv{"r"-(2,1)&"c"+(2,0)};
"t5";"t7" **i\crv{"t5"+(0,-10)&"t7"+(0,-10)} \POS?(.42)="s";
"t5";"s"*{\hole} **\crv{"t5"+(0,-5)&"s"-(5,0)};
"s"*{\hole};"t7" **\crv{"s"+(3,0)&"t7"+(0,-10)} \POS?(.5)*\dir{>}
\endxy
\]

\end{frame}

\begin{frame}
Freyd-Yetter (1989)  showed that $\mathbb{T}$ is an autonomous braided monoidal category with a ``freeness property''. 

\bigskip
We can see that $\mathbb{B}$ is the braided monoidal subcategory of $\mathbb{T}$ by identifying $n \in \mathbb{B}$ with  $\underbrace{+ \ldots +}_n \in \mathbb{T}$.   

\bigskip
The braids occur as the progressive endomorphisms in $\mathbb{T}$.   

\bigskip
The endomorphisms of the unit object 0 are precisely the \emph{oriented links}.

\end{frame}
\begin{frame}
\begin{proposition}
The trace of a braid in $\mathbb{T}$ is the link obtained as the Markov closure.
\end{proposition}
%
Here are examples of the trace of a braid, the second one is a trefoil.

\vskip-1.5cm\[
\xygraph{!{/r1.1pc/: /u1.2pc/::}
[uu] ( ?*+!D\txt\tiny{+}  !{+U\vcap[7]} [r(7)]="a7" , !{\vtwistneg}
( !{\xcapv-@(0)} []*+!U\txt\tiny{+} !{+D\vcap[-7]} [r(7)] - "a7"
 , [r] !{\vtwistneg} ( 
   ?*+!U\txt\tiny{+} !{+D\vcap[-5]} [r(5)] !{\POS= "a5"}
   , [r]*+!U\txt\tiny{+} !{+D\vcap[-3]}  [rrr]="a3"
 ) , [urrr] (
     ?*+!D\txt\tiny{+} !{+U\vcap} [r] !{\POS= "a1"}
    , [l]*+!D\txt\tiny{+} !{+U\vcap[3]} [rrr]- "a3"
    , [ll]*+!D\txt\tiny{+} !{+U\vcap[5]} [r(5)]- "a5"
 , [l]!{\vtwistneg}
   [r]!{\xcapv@(0)} []*+!U\txt\tiny{+} !{+D\vcap[-]} [r]- "a1"
)))}
\qquad\qquad\qquad\qquad
\xygraph{ !{0;/r1pc/:/u1.25pc/::} [uu]
 ( !{\vtwist\vtwist\vtwist}[uuu] 
   (?*+!D\txt\tiny{+} !{+U\vcap[3]}[rrr]-[d(4.45)] 
   , [r]*+!D\txt\tiny{+} !{+U\vcap}[r]-[d(4.45)]  )
 , [ddd]*+!U\txt\tiny{+} !{+D\vcap[-3]} 
 , [dddr]*+!U\txt\tiny{+} !{+D\vcap[-1]}
 )}
\]
\end{frame}
\begin{frame}
\frametitle{Invariants of links}

Each tensor-preserving (= strong monoidal) functor
\[
F : \mathbb{T}   \lra \mathrm{Vect}
\]
will produce a complex number for each oriented link $\ell$ :
\[
0 \stackrel{\ell}{\lra}  0   \longmapsto \mathbb{C}  \stackrel{F\ell}{\lra} \mathbb{C}\,.
\]
Such $F$ are determined by a single vector space $V$ and a special kind of Yang--Baxter operator
\[
y : V \otimes V \cong V \otimes V,
\]
on $V$ :
\[
F (+ + - + -) = V \otimes V \otimes V^\ast \otimes V \otimes V^\ast.
\]
There are techniques for solving the Yang-Baxter equation to find such functors, and hence, invariants of links; 
see Turaev (1988).
\end{frame}

\begin{frame}
\frametitle{Trace without duals}
A (\emph{right}) \emph{internal hom} $C^A$ for objects $A$ and $C$ of a monoidal category $\CC$ is an object equipped with a morphism (called ``evaluation'')
\[
\mathrm{ev} \ : A \otimes C^A \lra C
\]
which induces a bijection
\begin{align*}
\CC (B,C^A) &\cong \CC (A \otimes B, C)
\\
(B  \stackrel{f}{\lra} C^A) 
  &\longmapsto (A \otimes B \stackrel{1 \otimes f}{\lra} 
  A \otimes C^A \stackrel{\mathrm{ev}}{\lra}   C)\,. 
\end{align*}
If $A$ has a left dual $A^{\vee}$ then $C^A \cong A^{\vee}\otimes C$.  

Yet internal homs can exist without duals (e.g., in $\CC = \mathrm{Vect} $ and $ \mathrm{Set} $).

\end{frame}

\begin{frame}
\frametitle{Trace in analysis}

I know of no extension of the string notation, in keeping with the geometry as for duals, which covers internal homs.
\[
\xygraph{ !{/r2.5pc/: /u3pc/::}
 [d] *+[Fo][o]{g} (
  ? - [ul]^(.9)*{A} , ? - [ur]_(.9)*{B} , ? - [d]^(.9)*{C}
 )} \quad = \quad  
 \xygraph{ !{/r2.85pc/: /u2pc/::}
 [d(1.5)] *+[Fo][o]{\mathrm{ev}} (
  ? - [uul]^(.9)*{A} 
  , ? - [ur(.9)]*+[Fo][o]{\hat{g}} _*{C^A}
     - [r(.5)u(1.1)]_(.9)*{B}
  , ? - [d(1.2)]^(.9)*{C}
 )} 
\]

From field theory in physics there are geometric situations where internal homs and trace can exist without duals.   

This led Stolz-Teichner (arXiv:1010.4527) to develop a notion of \emph{trace} using only the ``weak duals'' $I^A$.

The following treatment of ``contraction'' is adapted from these authors. 

\end{frame}
\begin{frame}
\frametitle{Approximability and nuclearity}

A quadruple $(X,Y,A,B)$ of objects of a monoidal category $\CC$ is called \emph{approximable} when $Y^A$ exists and the following composite is injective.
\[
\xymatrix @C+2mm {
\CC (X,Y^A \otimes B) \ar[r]^-{A \otimes -}   
& \CC (A \otimes X, A \otimes Y^A \otimes B)  \ar[r]^-{\CC(1,\mathrm{ev} \otimes 1)} 
& \CC (A \otimes X, Y \otimes B)
}\]

A morphism $f : A \otimes X \lra Y \otimes B$ is called \emph{nuclear} when it is in the image of that composite.

\[
\xygraph{ !{/r2.85pc/: /u3pc/::}
 [d(1.5)] *+[Fo][o]{\mathrm{ev}} (
 ? - [l(.2)u(1.7)]^(.9)*{A} , ? - [l(.2)d(.8)]_(.9)*{Y} , 
 -  [u(.6)r] *+[Fo][o]{\hat{f}}_*{Y^A} (
  ? - [r(.2)d(1.4)]^(.9)*{B}, ? - [r(.3)u(1.1)]_(.9)*{X}
 ))} \quad = \quad  
 \xygraph{ !{/r2pc/: /u3pc/::}
 [d(1)] *+[Fo][o]{f} (
 ? - [ul]^(.9)*{A} , ? - [ur]_(.9)*{X} , ? - [dl]_(.9)*{Y} , ? - [dr]^(.9)*{B}
  )} 
\]
\end{frame}
\begin{frame}
\frametitle{More general trace}

\begin{example}
If $A$ has a right dual then any $(X,Y,A,B)$ is approximable. 
The appropriate composite is invertible so every $f$ is nuclear.
\end{example}

The \emph{trace} (or \emph{contraction by} $A$) of 
\[
f : A \otimes X \lra Y \otimes A
\]
in a braided monoidal  $\CC$ is the composite
\[
\mathrm{Tr} (f) : X \stackrel{\hat f}{\lra} Y^A \otimes A \stackrel{\gamma}{\lra} A \otimes Y^A \stackrel{\mathrm{ev}}{\lra} Y\,.
\]

\end{frame}

\begin{frame}

Let $\mathrm{Tv}$ be the category of topological complex vector spaces which are locally convex, complete and Hausdorff. 

\bigskip
The tensor product represents continuous bilinear maps.

\bigskip
\begin{theorem} [Stolz-Teichner]
Take $\CC = \mathrm{Tv}$.
%
\begin{enumerate}
\item If $A$ has the approximation property (from analysis) then $( \mathbb{C},  \mathbb{C}, A, A)$ is approximable.
\item A morphism $f:A\lra B$ is nuclear as in analysis if and only if $f:A \otimes \mathbb{C} \lra \mathbb{C} \otimes B$ is nuclear.
\item If $A$ has the approximation property and $f:A\lra A$ is nuclear then the trace $\mathrm{Tr}(f)$ is the classical trace.
\end{enumerate}
\end{theorem}
\end{frame}
\end{document}
